% Source: Wald GR, physical PDF pp. 196--200
%         (printed pp. 183--187).
% Coverage: Section 7.5, Perturbations; equations (7.5.1)--(7.5.23).
% Figures/tables/footnotes: no figures, tables, or footnotes.
% Status: source-reviewed and compile-clean.
% Modernization: metric perturbation gamma_ab -> h_ab (and trace reverse);
%                gravitational gauge terminology -> de Donder (harmonic);
%                electromagnetic “Lorentz gauge” -> Lorenz gauge.
% Uncertainties: none.

\subsection{Perturbations}\label{sec:7.5}

As emphasized at the beginning of this chapter, relatively few physically
interesting exact solutions of Einstein's equation are known. The known solutions,
such as the Robertson--Walker solutions (chapter 5) and the Schwarzschild solution
(chapter 6), may tell us a great deal about cosmology and the gravitational fields of
isolated bodies, but they leave many questions unanswered. For example, we might
wish to know how a small inhomogeneity in the matter distribution in a nearly
Robertson--Walker universe develops with time, or how a nearly spherical star
behaves if displaced slightly from equilibrium, or what happens if a small amount of
gravitational radiation is incident on a Schwarzschild black hole. It appears hopeless
to attempt to find exact solutions describing these processes. However, if the
deviation from a known exact solution is small, it makes sense to look for an
approximate solution by writing \(g_{ab}={}^{0}g_{ab}+h_{ab}\) (where
\({}^{0}g_{ab}\) is the known exact solution) and “linearizing” Einstein's equation
in \(h_{ab}\). We already followed this procedure in section 4.4 for the case where
\({}^{0}g_{ab}\) was the Minkowski metric, \(\eta_{ab}\). In this section, we shall
systematize the general procedure for deriving linearized equations, and for
Einstein's vacuum equation we shall obtain explicitly the vacuum perturbation
equations of an arbitrary exact solution.

Consider an equation
\begin{equation}
  \mathcal E(g)=0
  \tag{7.5.1}\label{eq:7.5.1}
\end{equation}
for an unknown function \(g\) (which, more generally, may be a collection of
functions or tensor fields, etc.). In the case of interest, \(g\) is the spacetime
metric possibly together with variables describing the matter distribution, and
\(\mathcal E\) is Einstein's equation. Suppose an exact solution, \({}^{0}g\), is
known and suppose we are interested in studying situations where the deviation
from \({}^{0}g\) is small. What we would really like to have is a one-parameter
(or multiparameter) family \(g(\lambda)\) of exact solutions,
\begin{equation}
  \mathcal E[g(\lambda)]=0 ,
  \tag{7.5.2}\label{eq:7.5.2}
\end{equation}
where \(\lambda\) measures the size of perturbation in the sense that
(i) \(g(\lambda)\) depends differentiably on \(\lambda\), and
(ii) \(g(0)={}^{0}g\). Thus, small \(\lambda\) corresponds to small deviations
from \({}^{0}g\), and a knowledge of \(g(\lambda)\) for small \(\lambda\) would
give us an exact perturbed solution. However, equation~\eqref{eq:7.5.2} may be too
difficult to solve. Nevertheless, we can derive a much simpler equation from
equation~\eqref{eq:7.5.2} by differentiating it with respect to \(\lambda\) and
setting \(\lambda\) equal to zero,
\begin{equation}
  \left.
    \frac{d}{d\lambda}\bigl[\mathcal E(g(\lambda))\bigr]
  \right|_{\lambda=0}=0 .
  \tag{7.5.3}\label{eq:7.5.3}
\end{equation}
Equation~\eqref{eq:7.5.3} is a linear equation for
\begin{equation}
  \gamma=\left.\frac{dg}{d\lambda}\right|_{\lambda=0} ;
  \tag{7.5.4}\label{eq:7.5.4}
\end{equation}
i.e., it can be expressed in the form
\begin{equation}
  \mathcal L(\gamma)=0 ,
  \tag{7.5.5}\label{eq:7.5.5}
\end{equation}
where \(\mathcal L\) is a linear operator. (Equation~\eqref{eq:7.5.5} is referred
to as the “linearization” of equation~\eqref{eq:7.5.1} about \({}^{0}g\).) Since
linear equations are generally much easier to solve than nonlinear equations, it may
be feasible to solve equation~\eqref{eq:7.5.5} even if
equation~\eqref{eq:7.5.1} is intractable. If we can solve
equation~\eqref{eq:7.5.5}, then \({}^{0}g+\lambda\gamma\) should yield a good
approximation to \(g(\lambda)\) for sufficiently small \(\lambda\), and issues of
physical interest thus can be investigated.

The above procedure provides a powerful tool for obtaining approximate solutions.
However, two important points should be kept in mind when employing perturbation
techniques: (1) It is, in general, very difficult to estimate the error involved in
replacing \(g(\lambda)\) by \({}^{0}g+\lambda\gamma\), i.e., to determine how
small \(\lambda\) must be in order that the approximate solution have sufficient
accuracy. (2) As derived above, existence of a one-parameter family of solutions
\(g(\lambda)\) implies the existence of a solution of equation~\eqref{eq:7.5.5},
where \(\gamma=(dg/d\lambda)|_{\lambda=0}\). However, the existence of a
solution \(\gamma\) of equation~\eqref{eq:7.5.5} does not necessarily imply the
existence of a corresponding one-parameter family of solutions, \(g(\lambda)\), of
\eqref{eq:7.5.2}, i.e., there may be spurious solutions of
equation~\eqref{eq:7.5.5}. Thus, the issue of “linearization stability,” i.e., the
existence of exact solutions corresponding to a solution of the linearized equations,
must be investigated before a perturbation analysis can be applied with complete
reliability.

Let us now derive the linearized vacuum Einstein equation~\eqref{eq:7.5.5} for a
metric perturbation \(h_{ab}\) of an exact solution \({}^{0}g_{ab}\) of Einstein's
equation in vacuum,
\begin{equation}
  {}^{0}R_{ab}=0 .
  \tag{7.5.6}\label{eq:7.5.6}
\end{equation}
To do so, we need to compute the Ricci tensor \(R_{ab}(\lambda)\) for the metric
\(g_{ab}(\lambda)\) in a useful form, specifically, in terms of quantities related to
the “background metric” \({}^{0}g_{ab}\). Differentiation of this expression with
respect to \(\lambda\) at \(\lambda=0\) then will yield the equation we seek.

Let \({}^{\lambda}\nabla_a\) denote the derivative operator associated with
\(g_{ab}(\lambda)\), and let \({}^{0}\nabla_a\) denote the derivative operator
associated with \({}^{0}g_{ab}\). According to the general analysis of section 3.1,
the difference between \({}^{\lambda}\nabla_a\) and \({}^{0}\nabla_a\) is
determined by a tensor field \(C^c{}_{ab}(\lambda)\), which, according to
equation~\eqref{eq:3.1.28}, is given by
\begin{equation}
  C^c{}_{ab}(\lambda)
    =\frac12g^{cd}(\lambda)
      \left\{
        {}^{0}\nabla_ag_{bd}(\lambda)
        +{}^{0}\nabla_bg_{ad}(\lambda)
        -{}^{0}\nabla_dg_{ab}(\lambda)
      \right\} .
  \tag{7.5.7}\label{eq:7.5.7}
\end{equation}
The Riemann tensor \(R_{abc}{}^d(\lambda)\) associated with \(g_{ab}(\lambda)\)
can be computed in terms of \({}^{0}R_{abc}{}^d\) and
\(C^c{}_{ab}(\lambda)\) by replacing the derivative operator
\({}^{\lambda}\nabla_a\) in the definition of the Riemann tensor,
equation~\eqref{eq:3.2.3}, by its expression in terms of \({}^{0}\nabla_a\) and
\(C^c{}_{ab}\). Proceeding in the same manner as in the derivation of
equation~\eqref{eq:3.4.3}, we find
\begin{equation}
  R_{abc}{}^d
    ={}^{0}R_{abc}{}^d
      -2{}^{0}\nabla_{[a}C^d{}_{b]c}
      +2C^e{}_{c[a}C^d{}_{b]e} .
  \tag{7.5.8}\label{eq:7.5.8}
\end{equation}
Thus, the Ricci tensor of \(g_{ab}(\lambda)\) is given by
\begin{equation}
  R_{ac}
    =-2{}^{0}\nabla_{[a}C^b{}_{b]c}
      +2C^e{}_{c[a}C^b{}_{b]e} ,
  \tag{7.5.9}\label{eq:7.5.9}
\end{equation}
where we have used the fact that \({}^{0}R_{ac}=0\) since \({}^{0}g_{ab}\) is a
solution of the vacuum Einstein equation.

Equation~\eqref{eq:7.5.9} expresses \(R_{ac}(\lambda)\) in a convenient form for
evaluating its derivative
\(\dot R_{ac}(\lambda)=(dR_{ac}/d\lambda)|_{\lambda=0}\), with respect to
\(\lambda\) at \(\lambda=0\). It follows immediately from its definition that
\(C^c{}_{ab}(\lambda)\) vanishes when \(\lambda=0\), so the term quadratic in
\(C^c{}_{ab}\) will make no contribution to \(\dot R_{ac}\). Thus, we obtain
\begin{equation}
  \dot R_{ac}
    =-2{}^{0}\nabla_{[a}\dot C^b{}_{b]c} ,
  \tag{7.5.10}\label{eq:7.5.10}
\end{equation}
where
\begin{equation}
  \dot C^c{}_{ab}
    =\left.\frac{dC^c{}_{ab}}{d\lambda}\right|_{\lambda=0} .
  \tag{7.5.11}\label{eq:7.5.11}
\end{equation}
From equation~\eqref{eq:7.5.7} we find that
\begin{equation}
  \dot C^c{}_{ab}
    =\frac12{}^{0}g^{cd}
      \left\{
        {}^{0}\nabla_ah_{bd}
        +{}^{0}\nabla_bh_{ad}
        -{}^{0}\nabla_dh_{ab}
      \right\} ,
  \tag{7.5.12}\label{eq:7.5.12}
\end{equation}
where
\begin{equation}
  h_{ab}
    =\left.\frac{dg_{ab}}{d\lambda}\right|_{\lambda=0}
  \tag{7.5.13}\label{eq:7.5.13}
\end{equation}
and we have used the fact that
\({}^{0}\nabla_a({}^{0}g_{bc})=0\). Thus, substituting
equation~\eqref{eq:7.5.12} in equation~\eqref{eq:7.5.10}, the linearized Einstein
equation for \(h_{ab}\) is found to be
\begin{align}
  0=\dot R_{ac}
    &=-\frac12{}^{0}g^{bd}
        {}^{0}\nabla_a{}^{0}\nabla_ch_{bd}
      -\frac12{}^{0}g^{bd}
        {}^{0}\nabla_b{}^{0}\nabla_dh_{ac} \notag\\
    &\quad+{}^{0}g^{bd}
        {}^{0}\nabla_b{}^{0}\nabla_{(c}h_{a)d} .
  \tag{7.5.14}\label{eq:7.5.14}
\end{align}
Since all quantities aside from \(h_{ab}\) now refer only to the background metric
\({}^{0}g_{ab}\), we shall in the following drop the superscript zero on the
background metric and its derivative operator, and we shall use the background
metric \(g_{ab}\) and its inverse \(g^{ab}\) to raise and lower indices. In this
notation, equation~\eqref{eq:7.5.14} becomes
\begin{equation}
  0=-\frac12\nabla_a\nabla_ch
    -\frac12\nabla^b\nabla_bh_{ac}
    +\nabla^b\nabla_{(c}h_{a)b} ,
  \tag{7.5.15}\label{eq:7.5.15}
\end{equation}
where \(h=h^a{}_a=g^{ab}h_{ab}\). Note that equation~\eqref{eq:7.5.15} agrees
with equation~\eqref{eq:4.4.4} in the special case of Minkowski spacetime, where
\(g_{ab}=\eta_{ab}\) and \(\nabla_a=\partial_a\).

Equation~\eqref{eq:7.5.15} can be simplified by a convenient choice of gauge. As
shown at the end of section C.2 of appendix C, \(h_{ab}\) and
\(h_{ab}+2\nabla_{(a}v_{b)}\) represent the same physical perturbation, where
\(v^a\) is an arbitrary vector field. As in the special case of perturbations of
Minkowski spacetime considered in section 4.4, we can solve the curved spacetime
wave equation (see theorem 10.1.2 of chapter 10),
\begin{equation}
  \nabla^b\nabla_bv_a+R_a{}^bv_b=-\nabla^b\overline h_{ab} ,
  \tag{7.5.16}\label{eq:7.5.16}
\end{equation}
where
\begin{equation}
  \overline h_{ab}\equiv h_{ab}-\frac12g_{ab}h ,
  \tag{7.5.17}\label{eq:7.5.17}
\end{equation}
and thereby impose the de Donder (harmonic) gauge condition
\begin{equation}
  \nabla^b\overline h_{ab}=0 .
  \tag{7.5.18}\label{eq:7.5.18}
\end{equation}
In this gauge, the trace of equation~\eqref{eq:7.5.15} yields
\begin{equation}
  \nabla_a\nabla^ah=0 .
  \tag{7.5.19}\label{eq:7.5.19}
\end{equation}
Hence, we can use the restricted gauge freedom
\(h_{ab}\longmapsto h_{ab}+2\nabla_{(a}w_{b)}\) with
\(\nabla^b\nabla_bw_a+R_a{}^bw_b=0\) to satisfy the curved space analogs of the
initial value conditions (4.4.34a) and (4.4.34b), thereby obtaining \(h=0\)
throughout the spacetime by theorem 10.1.2. Thus, for an arbitrary vacuum
perturbation \(h_{ab}\) of an arbitrary vacuum solution \(g_{ab}\), we can always
choose a \emph{transverse traceless gauge} whereby
\begin{align}
  \nabla^ah_{ab}&=0 ,
  \tag{7.5.20}\label{eq:7.5.20}\\
  h&=0 .
  \tag{7.5.21}\label{eq:7.5.21}
\end{align}
However, in general no analog of the radiation gauge conditions
\(h_{0\mu}=0\) (\(\mu=1,2,3\)) used in perturbations of flat spacetime can be
imposed in addition to \eqref{eq:7.5.20} and \eqref{eq:7.5.21}.

From the properties of the Riemann tensor, we have
\begin{align}
  \nabla^b\nabla_{(c}h_{a)b}
    &=\nabla_{(c}\nabla^bh_{a)b}
      +R^b{}_{(ca)}{}^dh_{db}
      +R^b{}_{(c|b|}{}^dh_{a)d} \notag\\
    &=\nabla_{(c}\nabla^bh_{a)b}
      +R^b{}_{ca}{}^dh_{db} ,
  \tag{7.5.22}\label{eq:7.5.22}
\end{align}
where in the second line we have used the fact that
\(R^b{}_{ca}{}^dh_{db}\) is symmetric in \(c\) and \(a\) and
\(R_c{}^d=R^b{}_{cb}{}^d=0\) since \(g_{ab}\) is a vacuum solution. Thus, in
the transverse traceless gauge, equations~\eqref{eq:7.5.20} and
\eqref{eq:7.5.21}, the linearized Einstein equation becomes simply
\begin{equation}
  \nabla^b\nabla_bh_{ac}-2R^b{}_{ac}{}^dh_{bd}=0 .
  \tag{7.5.23}\label{eq:7.5.23}
\end{equation}
This is remarkably similar in form to Maxwell's equation~\eqref{eq:4.3.15} in the
Lorenz gauge.

While the form of equation~\eqref{eq:7.5.23} is simple, it should be kept in mind
that, with the notable exception of a flat background metric, in practice
equation~\eqref{eq:7.5.23} comprises a very complicated system of coupled partial
differential equations. Success in perturbation analyses has been achieved only in a
few cases where the background metric has a great deal of symmetry or possesses
other simplifying properties, and even in these cases success has not usually been
achieved by a direct attack on equation~\eqref{eq:7.5.23} or
\eqref{eq:7.5.15}.

Finally, the issue of linearization stability---i.e., the existence of a one-parameter
family of exact solutions corresponding to a solution of the linearized equation---has
been studied extensively for the vacuum Einstein equation. If the background
spacetime \((M,g_{ab})\) is “closed”---i.e., if it possesses a compact, spacelike
Cauchy hypersurface \(\Sigma\) (see chapter 8)---then the Einstein equation is
linearization stable about \((M,g_{ab})\) if and only if \((M,g_{ab})\) does not
possess a Killing vector field. If \((M,g_{ab})\) possesses a Killing field, then it
is necessary that \(h_{ab}\) satisfy a second-order integral constraint involving the
Killing field in order that a one-parameter family \(g_{ab}(\lambda)\) exist. On the
other hand, linearization stability is believed to hold for asymptotically flat
perturbations of all asymptotically flat background spacetimes (even if Killing fields
are present), although this has been proven only for the flat background spacetime.
Details of the linearization stability analyses can be found in Fischer and Marsden
(1979).
